\documentclass[11pt, answers]{exam}
\usepackage{tikz}
\usetikzlibrary{decorations.pathreplacing}
\usepackage{amsmath}
\usepackage{amsthm}
\usepackage[boxed]{algorithm}
\usepackage[noend]{algpseudocode}
\usepackage{babel}
\usepackage{titlesec}
\usepackage{changepage}
\usepackage{amsfonts}
\usepackage{amssymb}
\usepackage{commath}
\usepackage{mathrsfs}
\usepackage{fancybox}
\usepackage{xcolor}
\usepackage{enumitem}
\usepackage[noend]{algpseudocode}
\usepackage[margin = 1in]{geometry}
\usepackage[utf8]{inputenc}
\usepackage{pgfplots}
\pgfplotsset{compat=1.18}
\usepackage{graphicx}
\usepackage{braket}
\usetikzlibrary{calc}
\newif\ifsteps
\stepstrue
%\stepsfalse

\renewcommand{\solutiontitle}{\noindent}
\renewcommand{\qedsymbol}{$\blacksquare$}
\newcommand{\x}{\times}
\newcommand{\ra}{\rightarrow}
\newcommand{\imp}{\implies}
\newcommand{\ex}{\exists}
\newcommand{\prob}{\textbf{P}}
\newcommand{\R}{\mathbb{R}}
\newcommand{\C}{\mathbb{C}}
\newcommand{\F}{\mathbb{F}}
\newcommand{\N}{\mathbb{N}}
\newcommand{\Z}{\mathbb{Z}}
\newcommand{\Q}{\mathbb{Q}}
\newcommand{\comp}{\mathsf{c}}
\newcommand{\topo}{\mathcal{T}}
\newcommand{\vter}{\mathit{v_{\text{ter}}}}
\newcommand{\vex}{\mathit{v_{\text{ex}}}}
\newcommand{\lang}{\mathcal{L}}
\newcommand{\ham}{\mathcal{H}}
\newcommand{\ess}{\mathcal{S}}
\newcommand{\avg}[1]{\langle {#1} \rangle}
\newcommand{\bigavg}[1]{\left\langle {#1} \right\rangle}
%\newcommand{\braket}[2]{\langle {#1} \mid {#2}\rangle}
%\newcommand{\ket}[1]{\left\lvert {#1}\right\rangle}
\newcommand{\allint}{\int_{-\infty}^{+\infty}}
\newcommand\todo[1]{\textcolor{red}{#1}}
\newcommand{\fullwidthseed}{\par\noindent\rule{\linewidth}{0pt}\vspace*{-\topskip}}
\newcommand{\vect}[1]{|{#1}\rangle}
\newcommand{\schro}{Schr{\"o}dinger}
\allowdisplaybreaks
\usepackage{calligra}
\DeclareMathAlphabet{\mathcalligra}{T1}{calligra}{m}{n}
\DeclareFontShape{T1}{calligra}{m}{n}{<->s*[2.2]callig15}{}
\newcommand{\scriptr}{\mathcalligra{r}\,}
\newcommand{\boldscriptr}{\pmb{\mathcalligra{r}}\,}
\usepackage{comment}
\includecomment{steps}

\begin{document}
	\begin{center}
		\Large \bf Physics 7501 Quiz 2 \\ 9/11/2026 \\ Sharon Ye
	\end{center}
	
	\section*{(1.2.1) Pauli matrices}
	
	What are the eigenvalues of the $2\times 2$ matrix
	$e^{i\theta\hat{Z}}$ where
	$\hat{Z}=\begin{pmatrix}1 & 0 \\ 0 & -1\end{pmatrix}$
	is the 3rd Pauli matrix?
	
	\begin{solution}
		For a diagonal matrix $M$, we exponentiate every element along the main diagonal to get 
		\begin{align*}
			e^M =
			\begin{pmatrix}
				e^{M_{11}} &  & {\text{\huge0}} \\
				\ & \ddots & \ \\
				\text{\huge0} & \ & e^{M_{nn}}
			\end{pmatrix}
		\end{align*}
		thus
		\begin{align*}
			e^{i\theta\hat{Z}} = 
			\begin{pmatrix}
				e^{i\theta} & 0 \\
				0 & e^{-i\theta}
			\end{pmatrix}
		\end{align*}
		Since this is a diagonal matrix, its eigenvalues are just the diagonal entries $e^{\pm i\theta}$
	\end{solution}
	
	\section*{(1.2.2) Linear algebra}
	
	Find the simultaneous (i.e.\ common) eigenkets of Pauli matrices
	$\sigma_x=\begin{pmatrix}0 & 1 \\ 1 & 0\end{pmatrix}$ and
	$\sigma_z=\begin{pmatrix}1 & 0 \\ 0 & -1\end{pmatrix}$.
	If you cannot find any, explain.
	
	\begin{solution}
		$\sigma_x$ and $\sigma_z$ have no simultaneous eigenkets. We know from homework that non-commuting Hermitian operators cannot share a common \textit{basis} of eigenkets. In this case however, we can make the stronger statement that they share \textit{none}. Suppose they shared a non-zero eigenket $\ket{\psi}$
		\begin{align*}
			[X,Z]|\psi\rangle
			&=(XZ-ZX)|\psi\rangle\\
			&=xz|\psi\rangle-zx|\psi\rangle\\
			&=0
		\end{align*}
		$\implies$
		\begin{align*}
			[X,Z]\ket{\psi}=\begin{pmatrix}0&-2\\2&0\end{pmatrix}
			\begin{pmatrix}
				a \\
				b
			\end{pmatrix}
			=
			\begin{pmatrix}
				-2a \\
				2b
			\end{pmatrix}
			\implies a = b = 0
		\end{align*}
		which contradicts our original presumption that $\ket{\psi}$ was non-zero.
	\end{solution}
	
	\section*{(1.2.3) Measurements}
	
	Consider the following normalized state of a single qubit:
	\begin{align*}
		|\psi\rangle
		&= \begin{pmatrix}
			\cos\theta\,e^{i\phi} \\
			\sin\theta
		\end{pmatrix}.
		\tag{1.3}
	\end{align*}
	
	\noindent For a measurement of Pauli operator
	$\hat{X}=\begin{pmatrix}0 & 1 \\ 1 & 0\end{pmatrix}$,
	what is the probability $\prob(X=+1)$ of measurement outcome $X=1$?
	Express it as a function of $\theta$ and $\phi$.
	
	\begin{solution}
		In the standard $\hat{Z}$-basis, the $X=+1$ state is given by
		\begin{align*}
			\ket{+} = \frac{1}{\sqrt{2}}
			\begin{pmatrix}
				1 \\
				1
			\end{pmatrix}
		\end{align*}
		The Born rule states that $\prob(X=+1)=|\braket{+|\,\psi\,}|^2$, thus
		\begin{align*}
			\braket{+|\,\psi\,} = \frac{1}{\sqrt{2}}\left(\cos\theta e^{i\phi} + \sin\theta\right)
		\end{align*}
		$\implies$
		\begin{align*}
			\prob(X=+1) = |\braket{+|\,\psi\,}|^2 &= \frac{1}{2}\left(\cos^2\phi+(e^{-i\phi}+e^{i\theta})\sin\theta\cos\theta + \sin^2\theta\right) \\
			&=\frac{1}{2}\left(1 + 2\sin\theta\cos\theta\cos\phi\right) \\
			&=\frac{1}{2}\left(1+\sin(2\theta)\cos\phi\right)
		\end{align*}
	\end{solution}
\end{document}
